\section{Data and Method}

\subsection{District Partisan Composition}

District-level presidential vote shares are derived from VEST 2020 precinct
returns~\citep{vest2022} areal-interpolated to district assignments via
\texttt{bisect label-analyze --types partisan}.
For enacted plans: interpolated to the 2022 enacted district boundaries.
For bisect plans: interpolated to the 2020 bisect district assignments.

\paragraph{Note on data availability.}
VEST precinct data for NC, WI, TX, FL, and PA (2020) were used in prior bisect
analysis (L-track partisan metrics, B.11 NC results).
The partisan analysis here uses the district-level outputs already computed
by the bisect pipeline rather than raw precinct data.

\subsection{Competitiveness Measures}

\begin{definition}[Competitive district]
A district is \emph{competitive} if the two-party presidential vote margin is
less than 10 percentage points: $|v_i^D - 0.50| < 0.10$, where $v_i^D$ is
the Democratic presidential vote share in district $i$.
\end{definition}

\begin{definition}[Efficiency gap]
Following Stephanopoulos \& McGhee~\citep{stephanopoulos2015partisan}:
$\mathrm{EG} = (\text{D wasted votes} - \text{R wasted votes}) / \text{total votes}$.
Negative EG indicates Democratic advantage (Republicans waste more votes);
positive indicates Republican advantage.
\end{definition}

\subsection{Plans Compared}

\begin{table}[h]
\centering\small
\caption{Plans compared for each study state.}
\label{tab:plans-compared}
\begin{tabular}{lll}
\toprule
State & Bisect plan & Enacted plan \\
\midrule
NC & \texttt{bisect state --state NC --year 2020} & NC 2022 enacted (post-Harper v.\ Hall) \\
WI & \texttt{bisect state --state WI --year 2020} & WI 2022 enacted \\
TX & \texttt{bisect state --state TX --year 2020} & TX 2021 enacted \\
FL & \texttt{bisect state --state FL --year 2020} & FL 2022 enacted \\
PA & \texttt{bisect state --state PA --year 2020} & PA 2022 enacted (Commonwealth Court) \\
\bottomrule
\end{tabular}
\end{table}
